\documentclass[11pt]{article}
\usepackage[a4paper,margin=0.82in]{geometry}
\usepackage[T1]{fontenc}
\usepackage[utf8]{inputenc}
\usepackage{lmodern,amsmath,amssymb,amsthm,microtype,booktabs,tabularx,xcolor}
\usepackage[numbers]{natbib}
\usepackage[colorlinks=true,linkcolor=blue!60!black,citecolor=blue!60!black]{hyperref}
\setlength{\emergencystretch}{4em}
\newtheorem{theorem}{Theorem}
\newtheorem{remark}{Remark}
\title{\textbf{The Dyadic Shadow Equals a Local Oscillation Obligation}}
\author{Liang Wang\\Huazhong University of Science and Technology}
\date{July 2026}
\begin{document}
\maketitle
\begin{abstract}
Every prefix inside a dyadic block is the left block-boundary prefix plus one local consecutive-block increment.  Consequently TPC-159 boundary control closes the full endpoint target exactly when a named-atom local oscillation bound is added.  No averaging or exceptional-endpoint deletion is permitted.
The classification is \texttt{REDUCTION\_\allowbreak{}L1}; no L2 arithmetic
progress or endpoint credit is claimed.
\end{abstract}

\section{Frozen target}
The selected route is \texttt{O161.bad\_\allowbreak{}endpoint\_\allowbreak{}pointwise\_\allowbreak{}fixed\_\allowbreak{}atom} \citep{TPC182}.  The required
signature is actual fixed-$h_0$ packet, named fixed atom, deterministic every
prefix, deterministic every scale, fixed-$X$ power at that atom, and actual
active support.  Throughout $h_0=2$ is only a source-backed data fact.

\section{Exact result}
\begin{theorem}
For a block boundary b and b<k<=b', |S\_\allowbreak{}k|<=|S\_\allowbreak{}b|+max\_\allowbreak{}{b<r<=b'}|S\_\allowbreak{}r-S\_\allowbreak{}b|.  Therefore the missing dyadic-shadow estimate is precisely a local maximal increment theorem at the same fixed atom.
\end{theorem}
\begin{proof}
For b<k<=b', S\_\allowbreak{}k=S\_\allowbreak{}b+(S\_\allowbreak{}k-S\_\allowbreak{}b), giving full-prefix control from boundary and local-increment control.  Conversely each increment is S\_\allowbreak{}k-S\_\allowbreak{}b, so it is at most twice a full-prefix maximum.  Thus the two maxima are equivalent up to factor two on the identical endpoint family; constant factors do not alter a fixed power.
\end{proof}

The smallest missing literal theorem is
\texttt{FIXED\_\allowbreak{}ATOM\_\allowbreak{}LOCAL\_\allowbreak{}MAXIMAL\_\allowbreak{}INCREMENT\_\allowbreak{}POWER\_\allowbreak{}SAVING}.

\section{Scoped stop and claim firewall}
No new scoped stop is declared in this paper.

The result is L0/L1 only.  Phase averages and Lebesgue-almost-every phase are
not named atoms.  Fixed $h_0=2$ is not decay.  Archive addressing is not
canonical physical representation, and empty-domain truth is not actual
active support.  The named-atom endpoint charge is zero, so the strict
$1/400$ budget is unpaid.  We claim no program-positive L2, prime-pair lower
bound, or twin-prime theorem.

\section{Reproducibility}
The adjacent JSON payload records the full six-axis signature, source lock,
typed result, exact scoped-stop cell, endpoint ledger, and claim firewall.
The audit artifact contains eight true mutation-rejection assertions.
\bibliographystyle{plainnat}
\bibliography{references}
\end{document}
